\chapter{Local Flatness, Covariant Differentiation, and Curvature}

\lectureoverview{A metric can always be made Euclidean through first
order at one chosen point. Curvature begins where that coordinate
freedom ends: in the second-order failure of neighboring local frames
to fit together. We develop that statement in three stages. First come
normal coordinates, then the covariant derivative and its connection,
and finally the noncommutativity of two covariant derivatives, which
defines the Riemann tensor and gives tidal gravity its invariant
meaning.}

\section{The Necessary Burden of Indices}

General relativity is difficult in a rather particular way. Its basic
ideas are not impossibly obscure, but a calculation quickly becomes
crowded with indices, derivatives, and symbols. More compressed
languages exist: vielbeins, differential forms, spinors, twistors, and
related devices can package long component equations into compact
objects. Learning those languages, however, is itself a substantial
task, and compact notation can conceal the machinery one most needs to
understand on a first encounter.

Maxwell's equations offer the model. Written component by component,
they amount to several separate equations. Vector notation collects
them into divergence and curl equations, making the structure easier
to see. General relativity admits analogous compressions, but the
component method remains the dependable route when an explicit
calculation must actually be done. We shall therefore keep the
indices visible. Once their grammar becomes familiar, the conceptual
path is much smoother, even though the arithmetic can still fill
pages.

Symbolic computer packages can help with that arithmetic. Given a
metric \(g_{mn}(x)\), such a package can construct Christoffel symbols,
Riemann tensors, Ricci tensors, and Einstein tensors. It can then
substitute these expressions into field equations. This does not
remove the need to understand what the objects mean. It merely
separates the mechanical expansion from the geometry.

\begin{classroomqa}
\audiencequestion{What were the active areas of research in general
relativity at the time of this lecture?}
\lecturerresponse{The lecture singled out several fronts. Quantum
gravity and the quantum physics of black holes were central problems
in fundamental theory. Classical astrophysics and numerical
relativity studied systems such as orbiting and merging black holes,
where solving Einstein's equations and predicting gravitational
radiation require much more than blindly feeding equations to a
computer. The same mathematical structures were also finding uses in
condensed-matter physics and fluid dynamics. The historical point was
equally striking: general relativity had moved from a comparatively
small specialty in the 1960s to one of the central languages of
theoretical physics, tightly connected with quantum field theory.
\lecturetimestamp{00:05:15}}
\end{classroomqa}

This is the last lecture devoted to Riemannian geometry by itself.
Beginning with the next lecture, the geometry will be interpreted as
gravity. The bridge is already visible. Distinguishing a genuine
gravitational field from an apparent field caused by accelerated or
curvilinear coordinates is mathematically the same problem as
distinguishing a curved metric from a flat metric written in awkward
coordinates.

\section{The Invariant Question}

A two-dimensional geometry is flat when Euclidean geometry works
intrinsically: straight lines, parallel lines, triangle relations, and
distance measurements all have their Euclidean form. It is tempting to
say that flatness means
\(g_{mn}=\delta_{mn}\), but that statement is coordinate dependent.
Even an ordinary plane has a position-dependent metric in polar
coordinates. The correct statement is that a flat region admits
coordinates in which the metric is constant and Euclidean throughout
that region.

Suppose instead that the metric is handed to us in some arbitrary
coordinates. One bad test would be to search through coordinate
systems one after another, transforming the metric each time and
hoping eventually to find a Cartesian form. There are infinitely many
possibilities, so this is not a practical criterion.

The right strategy is to construct a diagnostic from \(g_{mn}\) and
its derivatives. We want an object \(R\) with two decisive properties:
\[
\begin{aligned}
R=0
&\quad\Longleftrightarrow\quad
\text{locally flat geometry},\\
R\neq 0 \text{ at } P
&\quad\Longrightarrow\quad
\text{curvature at } P.
\end{aligned}
\]
Most importantly, \(R\) must be a tensor. A tensor that vanishes in
one coordinate system vanishes in every coordinate system, so no
coordinate search is needed.

\begin{figure}[H]
\centering
\includegraphics[width=\linewidth]{lecture_03_figure_01.png}
\caption{The desired invariant test: a metric need not visibly equal
\(\delta_{mn}\), while a vanishing curvature object should identify a
locally flat geometry.}
\lectureframe{00:18:20}
\end{figure}

The qualification \emph{locally} matters. A vanishing Riemann tensor
guarantees local Euclidean geometry; global identifications can still
produce nontrivial topology. A flat cylinder or flat torus, for
example, can be locally Euclidean without being globally the plane.

\section{Metric Geometry in the Small}

It helps to see what changes with dimension. A zero-dimensional space
is only a point. A connected one-dimensional Riemannian space can be
straightened locally into a line; globally it may be an interval, a
line, or a loop. A bug confined to a closed one-dimensional track can
measure its circumference, but bending the track in an ambient room
does not alter its intrinsic metric.

Two dimensions admit much richer intrinsic geometry: a plane, a
sphere, a lumpy terrestrial surface, and a torus need not agree even
locally, and their global topologies can differ as well. The metric is
the instrument that records their infinitesimal distances:
\begin{equation}
 ds^2=g_{mn}(x)\,dx^m dx^n.
\label{eq:l3-line-element}
\end{equation}
The quantity \(ds^2\) is a scalar. Consequently the changes in
\(dx^m\) under a coordinate transformation are exactly compensated by
the transformation of \(g_{mn}\).

\begin{figure}[H]
\centering
\includegraphics[width=\linewidth]{lecture_03_figure_02.png}
\caption{The invariant line element written with a position-dependent
metric.}
\lectureframe{00:19:00}
\end{figure}

The metric is nondegenerate and therefore has an inverse:
\begin{equation}
 g^{mn}g_{nr}=\delta^m{}_r.
\label{eq:l3-inverse-metric}
\end{equation}
The mixed Kronecker tensor \(\delta^m{}_r\) is the identity map. In
ordinary Riemannian geometry the metric is also positive definite,
\begin{equation}
 g_{mn}u^m u^n>0
\qquad\text{for every nonzero }u^m.
\label{eq:l3-positive}
\end{equation}
Lorentzian spacetime will later relax this positivity while retaining
nondegeneracy.

\begin{figure}[H]
\centering
\includegraphics[width=\linewidth]{lecture_03_figure_03.png}
\caption{The inverse metric and the identity tensor on the board.}
\lectureframe{00:20:40}
\end{figure}

The metric identifies vectors and covectors:
\begin{equation}
 V_m=g_{mn}V^n,
 \qquad
 V^m=g^{mn}V_n.
\label{eq:l3-raise-lower}
\end{equation}
Moving an index is therefore an operation performed by the metric, not
a typographical change.

\begin{figure}[H]
\centering
\includegraphics[width=\linewidth]{lecture_03_figure_04.png}
\caption{Raising and lowering a vector index with the metric and its
inverse.}
\lectureframe{00:24:40}
\end{figure}

If all components of \(g_{mn}\) are constant in a region, a constant
linear change of coordinates can diagonalize and normalize the metric
there. The geometry is then flat. The interesting information must
therefore be tied to how the metric varies with position.

\section{Normal Coordinates: Flat Through First Order}

Choose a point \(P\). The normal-coordinate theorem says that one can
construct coordinates centered at \(P\) for which
\begin{align}
 g_{mn}(P)&=\delta_{mn}, \label{eq:l3-normal-metric}\\
 \partial_r g_{mn}(P)&=0. \label{eq:l3-normal-first}
\end{align}
The lecture calls these Gaussian normal coordinates. Geometrically,
one draws axes through \(P\) and extends their coordinate lines as
straight as the geometry permits. The precise notion of straightness
will become geodesic motion.

\begin{figure}[H]
\centering
\includegraphics[width=\linewidth]{lecture_03_figure_05.png}
\caption{Normal coordinates chosen so that the metric is Euclidean at
the selected point.}
\lectureframe{00:28:50}
\end{figure}

There is no unique such frame. An orthogonal rotation at \(P\) preserves
\(g_{mn}(P)=\delta_{mn}\). What matters is not the orientation of the
axes but the existence of a frame in which both
\eqref{eq:l3-normal-metric} and \eqref{eq:l3-normal-first} hold.

\begin{classroomqa}
\audiencequestion{Do these conditions hold only at the selected point,
or throughout a neighborhood?}
\lecturerresponse{They hold at the selected point \(P\). One may center
a normal frame at any point, but one generally cannot use a single
coordinate system over a curved region and keep the metric Euclidean
with all first derivatives zero everywhere. Extending the same
coordinates away from \(P\) exposes higher derivatives of the metric.
\lecturetimestamp{00:29:53}}
\end{classroomqa}

At \(P\), the second derivatives are generally unavoidable:
\begin{equation}
 \partial_r\partial_s g_{mn}(P)\neq 0
 \qquad\text{in general}.
\label{eq:l3-normal-second}
\end{equation}
Thus the metric carries no invariant curvature information at zeroth
or first order at one point; coordinate freedom can remove it. The
first possible invariant information appears at second order.

\begin{figure}[H]
\centering
\includegraphics[width=\linewidth]{lecture_03_figure_06.png}
\caption{The normal-coordinate conditions: the metric is Euclidean,
its first derivatives vanish, and its second derivatives need not.}
\lectureframe{00:30:50}
\end{figure}

\subsection{Why the coordinate freedom is sufficient}

Take \(P\) to be the common origin of coordinates \(x^m\) and \(y^m\).
After using a linear transformation to arrange the metric at the
origin, consider the quadratic part of a further coordinate change:
\begin{equation}
 x^m
 =y^m+C^m{}_{nr}y^n y^r+O(y^3),
 \qquad C^m{}_{nr}=C^m{}_{rn}.
\label{eq:l3-quadratic-transform}
\end{equation}
In four dimensions the symmetric pair \(nr\) has
\[
\frac{4(4+1)}{2}=10
\]
independent choices. For each of the four values of \(m\), there are
ten coefficients \(C^m{}_{nr}\), giving forty quadratic coefficients.

The symmetric metric also has ten independent components in four
dimensions. Each component has four first derivatives, so the
conditions \(\partial_r g_{mn}(P)=0\) comprise forty equations. The
matching count shows that the quadratic coordinate freedom has exactly
the right size to impose the first-derivative conditions. The
normal-coordinate theorem supplies the existence statement; counting
alone is a useful structural check rather than a complete proof of
solvability for an arbitrary nonlinear system.

\begin{figure}[H]
\centering
\includegraphics[width=\linewidth]{lecture_03_figure_07.png}
\caption{The quadratic coordinate transformation and the ten
independent quadratic monomials in four dimensions.}
\lectureframe{00:33:50}
\end{figure}

\begin{classroomqa}
\audiencequestion{Is
\(x^m=y^m+C^m{}_{nr}y^n y^r+\cdots\) an attempted coordinate
transformation?}
\lecturerresponse{Yes. The coefficients \(C^m{}_{nr}\) are chosen so
that, in the new coordinates, the first derivatives of the metric
vanish at the common origin. The linear part keeps the Jacobian
invertible near that point, while the quadratic part adjusts the first
metric derivatives. \lecturetimestamp{00:35:41}}
\end{classroomqa}

\begin{classroomqa}
\audiencequestion{Is this what it means to say that a smooth geometry
is locally flat?}
\lecturerresponse{It is the relevant first-order statement. At a
sufficiently small scale a smooth surface is approximated by its
tangent plane, and normal coordinates express that fact through
\(g_{mn}(P)=\delta_{mn}\) and
\(\partial_r g_{mn}(P)=0\). Exact flatness on a neighborhood is
stronger: it also requires the curvature tensor to vanish there.
\lecturetimestamp{00:36:19}}
\end{classroomqa}

\begin{classroomqa}
\audiencequestion{Does the equality between forty coefficients and
forty equations really ensure that the equations can be solved?}
\lecturerresponse{The count identifies the required coordinate
freedom, while the theorem guarantees the local construction for a
smooth nondegenerate metric. The lecture proposes checking the
two-dimensional case directly as an exercise. The argument does not
depend on four dimensions; in \(d\) dimensions the same symmetry and
coordinate freedom fit together. \lecturetimestamp{00:36:45}}
\end{classroomqa}

\begin{classroomqa}
\audiencequestion{Why is there no constant term in the expansion, and
which term is first order?}
\lecturerresponse{The term \(y^m\) is the first-order term. A constant
term was omitted because the two coordinate origins were deliberately
chosen to represent the same point, so \(y=0\) implies \(x=0\). That
choice costs no generality. \lecturetimestamp{00:38:52}}
\end{classroomqa}

\section{Why Componentwise Differentiation Fails}

Now consider a covariant vector field \(V_m(x)\). In straight Cartesian
coordinates on flat space, a geometrically constant vector has
constant components. One might therefore try to define its derivative
as \(\partial_rV_m\). The difficulty appears as soon as the coordinate
grid bends.

\begin{figure}[H]
\centering
\includegraphics[width=\linewidth]{lecture_03_figure_08.png}
\caption{A curvilinear coordinate grid superimposed on a locally
straight frame.}
\lectureframe{00:40:10}
\end{figure}

Place the same geometric vector at neighboring points. Relative to a
fixed Cartesian frame the components do not change. Relative to a
curvilinear basis, however, the basis vectors themselves turn as one
moves, so the component list changes even though the vector does not.

\begin{figure}[H]
\centering
\includegraphics[width=\linewidth]{lecture_03_figure_09.png}
\caption{The same vector at neighboring points acquires different
components because the local coordinate basis has rotated.}
\lectureframe{00:42:30}
\end{figure}

The naive condition
\begin{equation}
 \partial_rV_m=0
\label{eq:l3-naive-derivative}
\end{equation}
can therefore be true in one coordinate system and false in another.
If \(\partial_rV_m\) were a rank-two tensor, that would be impossible:
vanishing in one frame would imply vanishing in every frame. Ordinary
partial derivatives of vector components are consequently not tensor
components.

\begin{figure}[H]
\centering
\includegraphics[width=\linewidth]{lecture_03_figure_10.png}
\caption{The deliberately naive zero-derivative condition beside the
changing curvilinear basis.}
\lectureframe{00:44:00}
\end{figure}

We need a derivative that removes the apparent change caused by the
basis while retaining the genuine change of the vector field.

\section{Defining the Covariant Derivative}

Differentiation at \(P\) only compares values in an arbitrarily small
neighborhood of \(P\). We may therefore begin in normal coordinates
centered there. Take a vector at a neighboring point, shift its tail
back to \(P\), and keep its normal-coordinate components fixed during
the shift. Comparing the transported vector with the vector already at
\(P\), then dividing by the separation, gives the ordinary derivative
in that locally straight frame.

\begin{figure}[H]
\centering
\includegraphics[width=\linewidth]{lecture_03_figure_11.png}
\caption{The local comparison: translate the neighboring vector in
normal coordinates before taking the difference.}
\lectureframe{00:46:50}
\end{figure}

Now express that geometrically defined derivative in arbitrary
coordinates. There are two contributions: the components \(V_m\) may
change, and the coordinate basis may change. Index structure fixes the
form of the correction:
\begin{equation}
 \nabla_rV_m
 =\partial_rV_m-\Gamma^t{}_{rm}V_t.
\label{eq:l3-covector-derivative}
\end{equation}
The quantities \(\Gamma^t{}_{rm}\) are the connection coefficients,
also called Christoffel symbols. The minus sign is appropriate for a
lower vector index. At the center \(P\) of normal coordinates,
\begin{equation}
 \Gamma^t{}_{rm}(P)=0,
\qquad
 \nabla_rV_m(P)=\partial_rV_m(P).
\label{eq:l3-gamma-normal-zero}
\end{equation}

\begin{figure}[H]
\centering
\includegraphics[width=\linewidth]{lecture_03_figure_12.png}
\caption{The covariant derivative of a covector: ordinary component
change minus the change induced by the moving basis.}
\lectureframe{00:49:50}
\end{figure}

\begin{classroomqa}
\audiencequestion{Does
\(\nabla_rV_m=\partial_rV_m-\Gamma^t{}_{rm}V_t\) hold in arbitrary
coordinates, and what does it become in normal coordinates?}
\lecturerresponse{The full formula holds in arbitrary coordinates. At
the selected point in normal coordinates the first metric derivatives,
and hence the Christoffel symbols, vanish, leaving only
\(\partial_rV_m\). Away from that point, even the same normal-coordinate
chart generally has nonzero connection coefficients.
\lecturetimestamp{00:50:22}}
\end{classroomqa}

The covariant derivative is so named because the entire object
\(\nabla_rV_m\) transforms covariantly as a tensor; the name is not
merely a comment about the lower derivative index. The Christoffel
term is precisely what repairs the nontensorial transformation of the
partial derivative.

\begin{figure}[H]
\centering
\includegraphics[width=\linewidth]{lecture_03_figure_13.png}
\caption{The complete covector derivative and its connection
correction.}
\lectureframe{00:51:30}
\end{figure}

\begin{classroomqa}
\audiencequestion{Is this related to the convective derivative used in
fluid flow?}
\lecturerresponse{The lecture answers cautiously that there is a
connection: fluid mechanics also distinguishes change carried by a
flow from change seen at a fixed coordinate point, and covariant
derivatives occur in fluid equations on curved coordinates or curved
spaces. The operators are not identical by definition, but both correct
a naive derivative for structure that is moving with the description.
\lecturetimestamp{00:52:31}}
\end{classroomqa}

For the torsion-free connection used here,
\begin{equation}
 \Gamma^t{}_{rm}=\Gamma^t{}_{mr}.
\label{eq:l3-torsion-free}
\end{equation}
More general geometries can possess torsion and need not obey this
symmetry. The present connection is the one relevant to the standard
metric formulation of general relativity.

\begin{figure}[H]
\centering
\includegraphics[width=\linewidth]{lecture_03_figure_14.png}
\caption{The Christoffel symbols identified as the symmetric
connection coefficients of the torsion-free geometry.}
\lectureframe{00:54:30}
\end{figure}

There is a useful physical analogy. Normal coordinates play the role of
a freely falling frame at one event: the coordinate correction can be
made to disappear there. Arbitrary curvilinear coordinates resemble an
accelerated elevator, where apparent forces enter because the frame
itself changes. The analogy is local, and curvature will measure why a
single freely falling frame cannot remove gravity throughout an
extended region.

\subsection{One correction for each tensor index}

For a covariant rank-two tensor \(T_{mn}\), each lower index receives
its own connection correction:
\begin{equation}
 \nabla_sT_{mn}
 =\partial_sT_{mn}
 -\Gamma^t{}_{sm}T_{tn}
 -\Gamma^t{}_{sn}T_{mt}.
\label{eq:l3-rank-two-derivative}
\end{equation}
To remember the rule, hold one index fixed as a spectator while
differentiating the other, and then repeat for the second index.
Contravariant indices receive corresponding terms with the opposite
sign.

\begin{figure}[H]
\centering
\includegraphics[width=\linewidth]{lecture_03_figure_15.png}
\caption{The covariant derivative of a rank-two covariant tensor, with
one Christoffel correction for each lower index.}
\lectureframe{01:01:50}
\end{figure}

\begin{classroomqa}
\audiencequestion{Does the symbol \(\partial_r\) itself refer only to
normal coordinates, and is the tensor formula restricted to that
frame?}
\lecturerresponse{No. The notation
\(\partial_r=\partial/\partial x^r\) is meaningful in any coordinate
system, and the full covariant-derivative formula is written for
arbitrary coordinates. Only at the chosen point of a normal frame do
the Christoffel corrections vanish. \lecturetimestamp{00:57:22}}
\end{classroomqa}

\begin{classroomqa}
\audiencequestion{When the covariant derivative reduces to the partial
derivative in normal coordinates, is that reduction only at the
selected point? Are the Christoffel symbols tensors?}
\lecturerresponse{Yes, the reduction is pointwise. The Christoffel
symbols are not tensors: they can vanish at \(P\) in normal coordinates
and be nonzero at the same point in another coordinate system. Their
nontensorial transformation is exactly what allows the combination in
\eqref{eq:l3-covector-derivative} to transform as a tensor.
\lecturetimestamp{00:58:13}}
\end{classroomqa}

\begin{classroomqa}
\audiencequestion{What will the covariant derivative actually be used
for?}
\lecturerresponse{For field equations. Gravitational fields vary from
point to point, so their equations must contain derivatives. At the
same time, the equations should have the same content in every
coordinate system. Covariant differentiation turns tensors into
tensors and therefore permits differential equations whose truth does
not depend on a peculiar reference frame. \lecturetimestamp{01:02:35}}
\end{classroomqa}

\begin{classroomqa}
\audiencequestion{Can the covariant derivative be multiplied by a
coordinate differential to describe the change between neighboring
points?}
\lecturerresponse{Yes. Along an infinitesimal displacement \(dx^s\),
the covariant change is \(dx^s\nabla_sT\). The partial term measures
the change of components, while the connection terms account for the
change of the coordinate basis. Even flat space in polar coordinates
requires that distinction. \lecturetimestamp{01:03:22}}
\end{classroomqa}

\section{Metric Compatibility Determines the Connection}

Apply the rank-two rule to the metric. At \(P\) in normal coordinates,
\(\partial_sg_{mn}=0\) and \(\Gamma^t{}_{sm}=0\), so
\(\nabla_sg_{mn}=0\). Since the covariant derivative is a tensor, this
pointwise result is true in every coordinate system:
\begin{equation}
 \nabla_sg_{mn}=0.
\label{eq:l3-metric-compatibility}
\end{equation}
This is metric compatibility. It says that parallel transport
preserves the metric, and hence preserves lengths and angles.

\begin{figure}[H]
\centering
\includegraphics[width=\linewidth]{lecture_03_figure_16.png}
\caption{The metric-compatibility equation begins the derivation of
the connection from the metric.}
\lectureframe{01:06:20}
\end{figure}

Expanding \eqref{eq:l3-metric-compatibility} gives
\begin{equation}
 0
 =\partial_sg_{mn}
 -\Gamma^t{}_{sm}g_{tn}
 -\Gamma^t{}_{sn}g_{mt}.
\label{eq:l3-compat-expanded}
\end{equation}
Write two more copies with cyclically permuted derivative and metric
indices:
\begin{align}
 0
 &=\partial_mg_{sn}
 -\Gamma^t{}_{ms}g_{tn}
 -\Gamma^t{}_{mn}g_{st},
\label{eq:l3-compat-m}\\
 0
 &=\partial_ng_{sm}
 -\Gamma^t{}_{ns}g_{tm}
 -\Gamma^t{}_{nm}g_{st}.
\label{eq:l3-compat-n}
\end{align}

\begin{figure}[H]
\centering
\includegraphics[width=\linewidth]{lecture_03_figure_17.png}
\caption{Three metric-compatibility equations with permuted indices.}
\lectureframe{01:10:20}
\end{figure}

Add \eqref{eq:l3-compat-m} and \eqref{eq:l3-compat-n},
then subtract \eqref{eq:l3-compat-expanded}. Using the torsion-free
symmetry \eqref{eq:l3-torsion-free}, the unwanted connection terms
cancel:
\begin{equation}
 \partial_mg_{sn}+\partial_ng_{sm}-\partial_sg_{mn}
 =2\Gamma^t{}_{mn}g_{st}.
\label{eq:l3-isolate-gamma-lowered}
\end{equation}

\begin{figure}[H]
\centering
\includegraphics[width=\linewidth]{lecture_03_figure_18.png}
\caption{Adding two compatibility equations and subtracting the third
cancels the unwanted Christoffel terms.}
\lectureframe{01:12:10}
\end{figure}

Multiplication by \(g^{rs}\) isolates the connection:
\begin{equation}
 \boxed{
 \Gamma^r{}_{mn}
 =\frac12 g^{rs}
 \left(
 \partial_mg_{sn}
 +\partial_ng_{sm}
 -\partial_sg_{mn}
 \right)}
\label{eq:l3-christoffel}
\end{equation}
Thus the Levi-Civita connection is built from the metric and its first
derivatives. This is the promised reason that the Christoffel symbols
vanish at the center of normal coordinates.

\begin{figure}[H]
\centering
\includegraphics[width=\linewidth]{lecture_03_figure_19.png}
\caption{The inverse metric isolates the Christoffel symbol from the
lowered connection equation.}
\lectureframe{01:16:10}
\end{figure}

\begin{classroomqa}
\audiencequestion{Why were all three metric-compatibility equations
set equal to zero?}
\lecturerresponse{At the selected point in normal coordinates, the
ordinary derivatives of the metric vanish and so do the Christoffel
symbols. Hence the covariant derivative of the metric is zero there.
To establish the statement at any other point, center a normal frame at
that point and repeat. Because \(\nabla_sg_{mn}\) is a tensor, its
vanishing is coordinate independent. \lecturetimestamp{01:14:27}}
\end{classroomqa}

\begin{classroomqa}
\audiencequestion{What are the dimensions of the metric and the
Christoffel array, and why are there forty independent Christoffel
components in four dimensions?}
\lecturerresponse{In \(d\) dimensions, \(g_{mn}\) is a symmetric
\(d\times d\) matrix. The connection has three indices. For each fixed
upper index \(r\), the symmetric lower pair \(mn\) has
\(d(d+1)/2\) choices. Thus a torsion-free connection has
\(d^2(d+1)/2\) independent components, which is forty when \(d=4\).
\lecturetimestamp{01:16:18}}
\end{classroomqa}

\begin{figure}[H]
\centering
\includegraphics[width=\linewidth]{lecture_03_figure_20.png}
\caption{The board during the component-count discussion: ten choices
for the symmetric lower pair and four choices for the upper index.}
\lectureframe{01:16:30}
\end{figure}

\begin{classroomqa}
\audiencequestion{Which metric appears in
\eqref{eq:l3-christoffel}? Must a different normal-coordinate metric
be used at every point?}
\lecturerresponse{The formula uses the components of the metric in
whatever global or local coordinates one has chosen for the
calculation. Normal coordinates are only the local device used to
derive the rule: move to a normal frame at one point, differentiate
there, transform the result back, and then discard the auxiliary frame.
Once \eqref{eq:l3-christoffel} is known, it is evaluated directly in
the arbitrary coordinates of interest. \lecturetimestamp{01:17:21}}
\end{classroomqa}

\begin{classroomqa}
\audiencequestion{Does either
\(\nabla_sg_{mn}=0\) or the Christoffel formula tell us that the space
is flat?}
\lecturerresponse{No. Metric compatibility holds for every
Levi-Civita connection, curved or flat. Christoffel symbols can be
nonzero even on a flat plane written in polar coordinates. If all
\(\Gamma^r{}_{mn}\) vanish throughout a region in one chart, that
chart is Cartesian there and the region is flat; but their nonvanishing
alone does not diagnose curvature. A tensor built from derivatives and
products of \(\Gamma\) is still needed. \lecturetimestamp{01:19:27}}
\end{classroomqa}

\section{A Cone Reveals What the Connection Cannot}

Consider a rounded cone. Away from its tip it can be approximated by an
ordinary cone made from paper. Cut a wedge from a flat disk and glue the
exposed edges together. Nothing in the remaining paper has been
stretched; locally, every point away from the join and tip is still
flat. The missing wedge is recorded by a deficit angle
\(\Delta\).

\begin{figure}[H]
\centering
\includegraphics[width=\linewidth]{lecture_03_figure_21.png}
\caption{The cone introduced as a geometry that is locally flat away
from its tip but has nontrivial curvature concentrated at the tip.}
\lectureframe{01:23:00}
\end{figure}

\begin{figure}[H]
\centering
\includegraphics[width=\linewidth]{lecture_03_figure_22.png}
\caption{Cutting a wedge from a disk and identifying its edges produces
the conical deficit angle.}
\lectureframe{01:24:20}
\end{figure}

The paper construction makes the distinction tangible. On the cut
sheet, transport a vector while keeping it parallel in the ordinary
plane. After the wedge edges are identified, a path that circles the
tip returns to its starting point with the transported vector rotated
relative to its initial direction.

\begin{figure}[H]
\centering
\includegraphics[width=0.78\linewidth]{lecture_03_figure_23.png}
\caption{The paper-cone demonstration: a flat sheet with a wedge
removed becomes a cone when its cut edges are joined.}
\lectureframe{01:25:10}
\end{figure}

\begin{figure}[H]
\centering
\includegraphics[width=\linewidth]{lecture_03_figure_24.png}
\caption{Vectors drawn on the cut wedge before the edges are identified.}
\lectureframe{01:27:40}
\end{figure}

This rotation is holonomy. It is not caused by a poor choice of
coordinates along one segment. Each small segment can be treated in a
locally flat frame. The effect appears only after those local
comparisons are composed around a closed loop. For the ideal cone the
net rotation equals the deficit angle, up to orientation.

\begin{figure}[H]
\centering
\includegraphics[width=\linewidth]{lecture_03_figure_25.png}
\caption{Parallel transport around the conical tip returns a vector
with a rotated direction: the holonomy records the enclosed curvature.}
\lectureframe{01:29:10}
\end{figure}

\section{Two Routes Around an Infinitesimal Loop}

Replace the finite loop by a tiny coordinate rectangle. Starting at
one corner, transport first in the \(r\) direction and then in the
\(s\) direction. Compare that result with transport first in \(s\) and
then in \(r\). In flat Cartesian coordinates the two routes agree. In
a curved geometry they generally do not.

\begin{figure}[H]
\centering
\includegraphics[width=\linewidth]{lecture_03_figure_26.png}
\caption{Two orderings of infinitesimal transport around a coordinate
rectangle.}
\lectureframe{01:31:40}
\end{figure}

The algebraic measure is the commutator of covariant derivatives:
\begin{equation}
 \left(\nabla_s\nabla_r-\nabla_r\nabla_s\right)V_n.
\label{eq:l3-commutator}
\end{equation}
Ordinary mixed partial derivatives commute on smooth functions. The
failure here comes from differentiating the connection terms as well
as the vector components.

\begin{figure}[H]
\centering
\includegraphics[width=\linewidth]{lecture_03_figure_27.png}
\caption{The infinitesimal loop represented by the difference between
the two orders of covariant differentiation.}
\lectureframe{01:33:20}
\end{figure}

To compute it, first regard
\[
 T_{rn}=\nabla_rV_n
\]
as a rank-two covariant tensor. Then use
\eqref{eq:l3-rank-two-derivative} to differentiate \(T_{rn}\) in the
\(s\) direction. Repeat with \(r\) and \(s\) interchanged and subtract.
The operation is lengthy but mechanical: expand, apply the product
rule, and cancel the terms containing second partial derivatives of
\(V_n\).

\begin{figure}[H]
\centering
\includegraphics[width=\linewidth]{lecture_03_figure_28.png}
\caption{Beginning the double covariant derivative by treating
\(\nabla_rV_n\) as a rank-two tensor.}
\lectureframe{01:35:40}
\end{figure}

Using the convention adopted on the board, the result is
\begin{equation}
 \left(\nabla_s\nabla_r-\nabla_r\nabla_s\right)V_n
 =R_{srn}{}^tV_t.
\label{eq:l3-riemann-commutator}
\end{equation}

\begin{figure}[H]
\centering
\includegraphics[width=\linewidth]{lecture_03_figure_29.png}
\caption{The commutator of covariant derivatives defines the action of
the curvature tensor on a covector.}
\lectureframe{01:38:30}
\end{figure}

For this ordering of the commutator,
\begin{equation}
\begin{split}
 R_{srn}{}^t
={}&
 \partial_r\Gamma^t{}_{sn}
 -\partial_s\Gamma^t{}_{rn}\\
&+\Gamma^p{}_{sn}\Gamma^t{}_{pr}
 -\Gamma^p{}_{rn}\Gamma^t{}_{ps}.
\end{split}
\label{eq:l3-riemann}
\end{equation}
Other texts may reverse the sign or arrange the indices differently;
the convention is fixed by writing both
\eqref{eq:l3-riemann-commutator} and \eqref{eq:l3-riemann}.

\begin{figure}[H]
\centering
\includegraphics[width=\linewidth]{lecture_03_figure_30.png}
\caption{The Riemann tensor: two derivatives of the connection and two
quadratic connection terms.}
\lectureframe{01:39:40}
\end{figure}

Equation \eqref{eq:l3-christoffel} shows why curvature first appears at
second order in the metric. Each \(\Gamma\) contains a first derivative
of \(g\). The derivative terms in \eqref{eq:l3-riemann} therefore
contain second derivatives of \(g\), while the quadratic terms contain
products of first derivatives. In normal coordinates at \(P\), the
connection itself vanishes but its derivatives generally do not, so
the curvature survives:
\[
\begin{aligned}
 \Gamma^t{}_{mn}(P)&=0,\\
 R_{srn}{}^t(P)
 &=\partial_r\Gamma^t{}_{sn}(P)
 -\partial_s\Gamma^t{}_{rn}(P).
\end{aligned}
\]
It is important to keep the derivative order straight: the Riemann
tensor uses first derivatives of \(\Gamma\), which are second order in
the metric.

This is the desired diagnostic. If every component of the Riemann
tensor vanishes throughout a region, that region is locally flat. If
any component is nonzero at a point, no change of coordinates can
remove the curvature there. A computer algebra system can test these
components from an analytic metric, but the concept is simply the
failure of parallel transport to close around an infinitesimal loop.

\section{Curvature as Tidal Gravity}

The geometric diagnostic has a direct physical meaning. Imagine a
flexible object made from linked triangles. In a flat region it can be
moved without stress. Push it into a region where the intrinsic metric
curves, and different parts of the linkage cannot all retain their
preferred separations and angles. The object must stretch, squeeze, or
shear.

\begin{figure}[H]
\centering
\includegraphics[width=\linewidth]{lecture_03_figure_31.png}
\caption{A linked object moved from a flat region into a curved bulge
must deform to follow the local metric.}
\lectureframe{01:44:10}
\end{figure}

That deformation is the geometric counterpart of tidal force. In a
gravitational field, neighboring freely falling particles can
accelerate relative to one another: one direction may be stretched
while another is squeezed. The Riemann tensor controls this relative
acceleration. In the language introduced later as geodesic deviation,
for tangent \(u^\mu\) and separation \(\xi^\mu\),
\begin{equation}
 \frac{D^2\xi^\mu}{D\tau^2}
 =-R^\mu{}_{\nu\rho\sigma}
 u^\nu\xi^\rho u^\sigma,
\label{eq:l3-geodesic-deviation-preview}
\end{equation}
with the overall sign following the chosen curvature convention. This
equation is not derived here, but it makes precise the lecture's
statement that curvature and tidal stress are literally the same local
physics.

An ideal perfectly uniform gravitational field has no tidal variation
and can be removed locally by passing to a uniformly accelerated frame.
Real fields produced by isolated masses are not perfectly uniform.
Earth's field points approximately radially and changes with position,
so neighboring free-fall trajectories converge or separate.

\begin{figure}[H]
\centering
\includegraphics[width=\linewidth]{lecture_03_figure_32.png}
\caption{A radial gravitational field varies across an extended body,
producing the relative acceleration identified as tidal gravity.}
\lectureframe{01:48:40}
\end{figure}

\begin{classroomqa}
\audiencequestion{What distinguishes a genuine gravitational field
from the apparent field seen in an accelerated frame?}
\lecturerresponse{A uniformly accelerated frame can imitate a uniform
gravitational field for local experiments, but it does not produce
tidal relative acceleration. A field generated by real matter varies
from place to place and therefore produces tidal forces. Those forces
are invariant evidence that the effect is not merely a coordinate
transformation. A perfectly uniform real field is an idealization;
near a sufficiently small patch of a large body's surface, the actual
field can only be approximately uniform. \lecturetimestamp{01:46:54}}
\end{classroomqa}

\begin{classroomqa}
\audiencequestion{Why does the size of the falling object matter for
tidal effects?}
\lecturerresponse{Curvature compares neighboring locations. An object
much smaller than the length scale over which the field changes
samples almost the same acceleration at all of its parts and suffers
little deformation. An extended object samples a larger difference.
A freely falling bacterium near Earth feels negligible tidal stress,
whereas the lecture's imaginary two-thousand-mile person would be
severely distorted. The curvature has a local value, but the accumulated
deformation scales with the separation being compared.
\lecturetimestamp{01:49:01}}
\end{classroomqa}

\begin{classroomqa}
\audiencequestion{If flatness requires checking the curvature at every
point, is the test still an infinite problem?}
\lecturerresponse{For a metric given analytically, the derivatives and
curvature components are analytic expressions. One calculates them
once as functions of position and can often see whether they vanish
identically. If the metric is supplied only as a numerical table, then
one must estimate and inspect curvature throughout the sampled region,
which can indeed be expensive. Even then, this is far better defined
than searching every possible coordinate system at every point.
\lecturetimestamp{01:50:37}}
\end{classroomqa}

\section{What Has Been Established}

The argument can now be read as one continuous chain:
\begin{enumerate}
 \item At any chosen point, coordinates can make
 \(g_{mn}=\delta_{mn}\) and \(\partial_rg_{mn}=0\).
 \item Ordinary derivatives of tensor components fail to transform as
 tensors because coordinate bases vary from point to point.
 \item The connection corrects that basis variation and defines
 covariant differentiation.
 \item Metric compatibility and zero torsion determine the connection
 uniquely from the metric.
 \item Two covariant derivatives fail to commute by the Riemann tensor.
 That failure measures holonomy, second-order metric variation, and
 tidal gravity.
\end{enumerate}

The promised curvature diagnostic has therefore been constructed
without choosing a privileged global coordinate system. The next step
is to change from positive-definite Riemannian geometry to Lorentzian
spacetime, apply these tools to the Schwarzschild geometry of a star or
black hole, and return to the geodesics that were deferred at the end
of this lecture.
